\documentclass[a4paper,10pt]{article}
\usepackage[bottom=0.5in]{geometry}
%A Few Useful Packages
\usepackage{marvosym}
\usepackage{fontspec} 					%for loading fonts
\usepackage{xunicode,xltxtra,url,parskip} 	%other packages for formatting
\RequirePackage{color,graphicx}
\usepackage[usenames,dvipsnames]{xcolor}
\definecolor{date}{HTML}{666666} 
\definecolor{primary}{HTML}{2b2b2b} 
\definecolor{headings}{HTML}{6A6A6A}
\definecolor{subheadings}{HTML}{333333}

\usepackage[big]{layaureo} 				%better formatting of the A4 page
% an alternative to Layaureo can be ** \usepackage{fullpage} **
\usepackage{supertabular} 				%for Grades
\usepackage{titlesec}					%custom \section

%Setup hyperref package, and colours for links
\usepackage{hyperref}
\definecolor{linkcolour}{rgb}{0,0.2,0.6}
\hypersetup{colorlinks,breaklinks,urlcolor=linkcolour, linkcolor=linkcolour}

%FONTS
\defaultfontfeatures{Mapping=tex-text}
%\setmainfont[SmallCapsFont = Fontin SmallCaps]{Fontin}
%%% modified for Karol Kozioł for ShareLaTeX use

\setmainfont[Color=primary, Path = fonts/lato/]{Lato-Lig}[SmallCapsFont = Lato-Reg.ttf,BoldFont = Lato-Bol.ttf,ItalicFont = Lato-RegIta.ttf]
\setsansfont[Scale=MatchLowercase,Mapping=tex-text, Path = fonts/raleway/]{Raleway-ExtraLight}
\newcommand{\custombold}[1]{\color{subheadings}\fontspec[Path = fonts/lato/]{Lato-Reg}\selectfont #1 \normalfont}

%%%

%CV Sections inspired by: 
%http://stefano.italians.nl/archives/26
\titleformat{\section}{\Large\scshape\raggedright}{}{0em}{}[\titlerule]
\titlespacing{\section}{0pt}{3pt}{3pt}
%Tweak a bit the top margin
%\addtolength{\voffset}{-1.3cm}

%Italian hyphenation for the word: ''corporations''
\hyphenation{im-pre-se}

%-------------WATERMARK TEST [**not part of a CV**]---------------
\usepackage[absolute]{textpos}
\usepackage[document]{ragged2e}
\setlength{\TPHorizModule}{30mm}
\setlength{\TPVertModule}{\TPHorizModule}
\textblockorigin{2mm}{0.65\paperheight}
\setlength{\parindent}{0pt}

\newcommand{\namesection}[3]{
	\centering{
		\sffamily
		\fontspec[Path = fonts/lato/]{Lato-Lig}\fontsize{40pt}{10cm}\selectfont #1 
		\fontspec[Path = fonts/lato/]{Lato-Reg}\selectfont #2
	} \\
	\vspace{5pt}
	\centering{
	    
	    \fontspec[Path = fonts/raleway/]{Raleway-Regular}
	    \fontsize{11pt}{14pt}\selectfont #3
	}
	\vspace{5pt}
}

%--------------------BEGIN DOCUMENT----------------------
\begin{document}

%WATERMARK TEST [**not part of a CV**]---------------
%\font\wm=''Baskerville:color=787878'' at 8pt
%\font\wmweb=''Baskerville:color=FF1493'' at 8pt
%{\wm 
%	\begin{textblock}{1}(0,0)
%		\rotatebox{-90}{\parbox{500mm}{
%			Typeset by Alessandro Plasmati with \XeTeX\  \today\ for 
%			{\wmweb \href{http://www.aleplasmati.comuv.com}{aleplasmati.comuv.com}}
%		}
%	}
%	\end{textblock}
%}

\pagestyle{empty} % non-numbered pages

\font\fb=''[cmr10]'' %for use with \LaTeX command

%--------------------TITLE-------------
\namesection{Calvin}{Xu}{ 
\href{mailto:cx23@illinois.edu}{\custombold{cx23@illinois.edu}}\\
\href{https://github.com/cmxu}{\custombold{github://cmxu}} | \href{https://cmxu.io}{\custombold{cmxu.io}} | \href{https://stackoverflow.com/users/5037048/cmxu}{\custombold{stack://cmxu}}
}

\section{Research Statement}
\justify{I am focused on making trustworthy ML techniques \textit{usable} in the \textit{real world}. My research lies at the intersection between \textit{ML} and \textit{formal methods}. I am broadly interested in the generation of practical adversarial examples, certified training, and NN verification for the vision, language, and wireless. I am currently working on LLM reasoning, agent synthesis,  and coding performance/safety.}

% I study the landscape adversarial examples, towards verifiable robustness while maintaining high accuracy. I am also interested in the generation of practical adversarial examples, robust training, neural network verification, and the intersection of explainable/interpretable with robust ML.

\section{Education}
\begin{tabular}{r|p{11cm}}
 \textsc{Current} & \textsc{PhD Student} at \textsc{University of Illinois at Urbana-Champaign}\\
 \textsc{Sep 2021} & \small{Working on \textbf{PhD} in \textbf{CS}, GPA: 4.0/4.0}\\\multicolumn{2}{c}{} \\
 \textsc{Dec 2019} & \textsc{Advanced Studies Student} at \textsc{Massachusetts Institute of Technology} \\
 \textsc{Aug 2019} & \small{Coursework - Advances in Computer Vision (6.869), GPA: 5.0/5.0}\\\multicolumn{2}{c}{} \\
 \textsc{Dec 2018} & \textsc{Undergraduate/Graduate Student} at \textsc{Washington University in St. Louis} \\
 \textsc{Aug 2015} & \small{\textbf{MS} in \textbf{CS}, Certificate in Data Mining \& Machine Learning, GPA: 3.9/4.0}\\
 & \small{\textbf{BS} in \textbf{Mathematics} and \textbf{BS} in \textbf{CS}, GPA: 3.8/4.0}\\
\end{tabular}

\section{Publications}
\textbf{Compression Aware Certified Training}\\
\underline{Changming Xu}, Gagandeep Singh. Arxiv Preprint\\[-8pt]

\textbf{Misaligning Reasoning with Answers--A Framework for Assessing LLM CoT Robustness}\\
Enyi Jiang, \underline{Changming Xu}, Nischay Singh, Gagandeep Singh. Arxiv Preprint\\[-8pt]

\textbf{Safety and Trust in Artificial Intelligence with Abstract Interpretation}\\
Gagandeep Singh, Jacob Laurel, Sasa Misailovic, Debangshu Banerjee, Avaljot Singh, \underline{Changming Xu}, Shubham Ugare, Huan Zhang. FTPL\\[-8pt]

\textbf{Support is all you need for Certified VAE Training}\\
\underline{Changming Xu}, Debangshu Banerjee, Deepak Vasisht, Gagandeep Singh. ICLR '25\\[-8pt]

\textbf{Certified DNN Training against Universal Adversarial Perturbations}\\
\underline{Changming Xu}, Gagandeep Singh. ECCV '24\\[-8pt]

\textbf{Scalable Relational Verification and Training for Deep Neural Networks}\\
Debangshu Banerjee, \underline{Changming Xu}, Gagandeep Singh. SAIV@CAV '24\\[-8pt]

\textbf{Robust Universal Adversarial Perturbations}\\
\underline{Changming Xu}, Gagandeep Singh. ICML '24\\[-8pt]

\textbf{Input-Relational Verification of Deep Neural Networks}\\
Debangshu Banerjee, \underline{Changming Xu}, Gagandeep Singh. PLDI '24\\[-8pt]

\textbf{Bypassing the Safety Training of Open-Source LLMs with Priming Attacks}\\
Jason Vega*, Isha Chaudhary*, \underline{Changming Xu}*, Gagandeep Singh. ICLR Tiny Paper '24 \\[-8pt]

\textbf{Exploring Practical Vulnerabilities of Machine Learning-based Wireless Systems}\\
Zikun Liu, \underline{Changming Xu}, Emerson Sie, Deepak Vasisht, Gagandeep Singh. NSDI '23\\[-8pt]

\textbf{Race Detection and Reachability in Nearly Series-Parallel DAGs}\\
Kunal Agrawal, Joseph Devietti, Jeremy T. Fineman, I-Ting Angelina Lee, Robert Utterback, \underline{Changming Xu}. SODA '18\\
\hspace*{6pt}{\footnotesize $^*$equal contribution}%, $^\dagger$in submission}

\section{Fellowships \& Awards}
\begin{tabular}{r|p{11cm}}
\textsc{Aug} 2022 & \href{https://www.qualcomm.com/research/university-relations/innovation-fellowship/2022-north-america}{Qualcomm Innovation Fellowship Winner} (one of 19 teams out of $\sim$300)\\
& \textit{Proposal}: Provably Robust Machine Learning for Wireless Systems \\
\textsc{Aug} 2019 & 1st place at \href{https://sites.google.com/view/advml/advml19-challenge}{AdvML Challenge} at \textit{SigKDD 2019} \\
\textsc{Apr} 2018 & \href{https://eece.wustl.edu/academics/graduate-programs/Graduate-Fellowships-Scholarships.html}{Dean's Select Fellowship for Research Excellence (WUSTL)} \\
\textsc{Dec} 2015-17 & Putnam Exam: 28, 20, 12 \\
\textsc{Aug} 2015 & Compton Scholar for Mathematics and Physics (4 per grade) \\
\textsc{Apr} 2017 & Missouri Math Competition: 1st Place Team \\
\textsc{Dec} 2015 & ICPC Regional: Top 5 Team
\end{tabular}


% \begin{tabular}{l|l}
% Robust Universal Adversarial Perturbations* & ICML '24\\
% Exploring Practical Vulnerabilities of Machine Learning-based Wireless Systems & NSDI '23\\
% \href{https://dl.acm.org/doi/10.5555/3174304.3175277}{Race Detection and Reachability in Nearly Series-Parallel DAGs} & SODA '18\\
% \end{tabular}\\


\section{Work Experience}
\begin{tabular}{r|p{11cm}}
\textsc{Sept 2025} & \textsc{Intern} at \textsc{Meta} \\
 \textsc{May 2025} & \small{\vspace{-10pt}\begin{itemize}
 \item LLM data generation and classification for notification ranking model \end{itemize}\vspace{-15pt}}\\\multicolumn{2}{c}{} \\
\textsc{Dec 2024} & \textsc{Intern} at \textsc{Bytedance} \\
 \textsc{Sept 2024} & \small{\vspace{-10pt}\begin{itemize}
 \item Finetuning and soft prompting VLMs for content moderation \end{itemize}\vspace{-15pt}}\\\multicolumn{2}{c}{} \\
 \textsc{May 2021} & \textsc{Associate Staff} at \textsc{MIT Lincoln Laboratory} \\
 \textsc{Mar 2019} & \small{\vspace{-10pt}\begin{itemize}
 \item Applied DAGAN to network traffic to augment unbalanced classes resulting in a \textasciitilde 20\% increase in classification accuracy
 \item Applied CNNs, Word2Vec, and Q-Learning techniques to flight data to predict flight reroutes, published at \textit{INFORMS '19}
 \item Gathered dataset and built system using NLP and Neural Networks (NNs) to determine when to apply Adversarial ML
 \item Explored Semantic Adversarial Examples and Spatial Transformer Networks for data augmentation
 \item Developed an algorithm using GANs and Non-Negative NNs that won 1st place at \href{https://sites.google.com/view/advml/advml19-challenge}{AdvML Challenge} at \textit{SigKDD 2019}\end{itemize}\vspace{-15pt}}\\\multicolumn{2}{c}{} \\
 \textsc{Dec 2018} & \textsc{Teaching Assistant} at \textsc{Washington University in St. Louis}\\
 \textsc{Sep 2016} & \small{\vspace{-10pt}\begin{itemize}
 \item Graded and held office hours for graduate level Machine Learning (150 students) and Numerical Applied Mathematics (50 students)
 \item Designed, ran, and graded the final project for the Machine Learning course\end{itemize}\vspace{-15pt}}\\\multicolumn{2}{c}{} \\
 \textsc{Aug 2018} & \textsc{Intern} at \textsc{EUB-INC} in Beijing, China\\
 \textsc{May 2018} & \small{\vspace{-10pt}\begin{itemize}
 \item Leveraged NNs and clustering to automate user grouping for data-driven advertisement on WeChat. 
 \item Developed algorithm which reduced turnaround time by 90\% for data team, allowing for faster, more targeted advertisement\end{itemize}\vspace{-15pt}}\\
\end{tabular}

\section{Research Experience}
\begin{tabular}{r|p{11cm}}
 \textsc{Current} & \textsc{PhD Student} at \textsc{University of Illinois in Urbana-Champaign}\\
 \textsc{Aug 2021} & \small{Currently working on 
 \begin{itemize}
     \item  Certified training and efficient DNN verification of universal adversarial perturbations, VAEs, and other attacks/architectures
     \item Training for better LLM alignment
     \item Training networks that are certifiably robust under network compression (pruning/quantization)
     \item Probabilistic verification of VAEs in the wireless domain
     \item Certifiably robust training of DNNs
 \end{itemize}\vspace{-15pt}}\\\multicolumn{2}{c}{} \\
 % \small{Developed technique for generating Robust UAPs which have strong potential for creating real-world attacks, submitted to \textit{ICML '23}. Created and implemented technique for unsupervised UAP generation. Working on attacks and verification techniques for wireless systems, accepted at \textit{NSDI '23}. Currently focused on studying the landscape of adversarial examples and robustness against UAPs.}\\\multicolumn{2}{c}{} \\
 \textsc{Mar 2018} & \textsc{Research Assistant} at \textsc{Carnegie Mellon University}\\
 \textsc{Jun 2017} & \small{\vspace{-10pt}\begin{itemize}
     \item Employed Adversarial ML techniques to thwart defect prediction
     \item Developed theory for attacking and defending high dimensional SVMs
     \item Implemented data poisoning attacks on the \href{https://www.sec.cs.tu-bs.de/~danarp/drebin/}{Drebin Android Malware} dataset, and presented poster at CMU
     \item Proved that \textasciitilde 10 poisoned data points can be enough to significantly reduce the effectiveness of the malware detector
 \end{itemize}\vspace{-15pt}}\\\multicolumn{2}{c}{} \\
 \textsc{May 2017} & \textsc{Research Assistant} at \textsc{Washington University in St. Louis}\\
 \textsc{Jun 2016} & \small{\vspace{-10pt}\begin{itemize}
 \item Created benchmarks and proved correctness for a work-stealing scheduler, improving cache efficiency by factor of 2
 \item Rigorously proved a nearly series parallel race detection algorithm which matched asymptotic runtime of series parallel case
 \item Published paper in \textit{ACM-SIAM SODA '18}\end{itemize}\vspace{-15pt}}\\
\end{tabular}

% %Section: Languages
% \section{Skills}
% \begin{tabular}{r|p{11cm}}
% \textsc{Languages} & English (Native), Chinese (Fluent), Japanese (Basic Knowledge)\\
% \end{tabular}

\iffalse
\section{Interests and Activities}
Math Club (President)\\
Bridge Club (Founder)\\
Chess Club\\
Badminton\\
\fi

%\section{Coursework}
%\begin{tabular}{r|p{11cm}}
% \textsc{Math} & Differential Equations, Matrix Algebra, Calculus of Several Variables, Introduction to Analysis, Number Theory and Cryptography, Numerical Applied Mathematics, Introduction to Lebesgue Integration, Topics in Applied Mathematics: Mathematics for Multimedia, Linear Algebra, Probability, Bayesian Statistics, Mathematical Statistics, Time Series Analysis, Statistical Computation \\\multicolumn{2}{c}{} \\
%\textsc{CS} & Programming Skills Workshop, Algorithms and Data Structures, \textit{Machine Learning}, Object Oriented Software Development Laboratory, Introduction to System Software, \textit{Introduction to Artificial Intelligence}, \textit{Computational Geometry}, \textit{Data Mining}, \textit{Multi-Agent Systems}, \textit{Topological Applications for Data Analysis and Machine Learning}, \textit{Approximation Algorithms}, \textit{Algorithms for Nonlinear Optimization}, \textit{Sparse Modeling for Imaging and Vision}, \textit{Systems Security}, \textit{Algorithms for Computational Biology}, \textit{Advanced Topics in Computer Vision}$^\dagger$, \textit{Statistical Reinforcement Learning}', \textit{Logic and AI}', \textit{Adv Topics in Sec, Priv and ML}'\\
%\end{tabular}\\
%\hspace*{37pt}\footnotesize{\textit{Graduate Level}, $^\dagger$ MIT, ' UIUC}

\end{document}
